\section{Proof of~\Cref{thm:main-formal}}
The purpose of this section is to prove our direct sum theorem (\Cref{thm:main-formal}) showing that if $f$ is hard to compute, than $f^{\otimes k}$ is even harder to compute. We will in fact prove a {\sl threshold} direct sum theorem, showing that it is hard even to get {\sl most} of the $k$ copies correct. To formalize this, we generalize the notation $\odepth(\cdot, \cdot)$ to take in an additional threshold parameter $t$ specifying how many blocks we allow to be wrong.
\begin{definition}
    \label{def:threshold-DPT}
    For any function $f:\bn \to \bits$, error $\eps$, and threshold $t \in \N$, we use $\odepthmuk(f^{\otimes}, \eps, t)$ to denote the minimum expected depth of a tree $T:(\bn)^k \to \bits^k$ satisfying
    \begin{equation*}
        \Prx_{\bX \sim \mu^k}\bracket*{\lzero{T(\bX) - f^{\otimes k}(\bX)}] > t} \leq \eps.
    \end{equation*}
\end{definition}

\begin{theorem}[A strong threshold direct sum theorem for query complexity]
    \label{thm:threshold-DPT}
    For every function $f:\bits^n \to \bits$, distribution $\mu$ over $\bits^n$, $k \in \N$, and $\gamma, \delta \in (0,1)$,
    \begin{equation*}
        \odepthmuk(f^{\otimes k}, 1 - e^{-\Omega(\delta k)} - \gamma, \Omega(\delta k)) \geq \Omega\left( \frac{\gamma^2 k}{\log(1/\delta)}\right) \cdot \odepthmu(f,\delta).
    \end{equation*}
\end{theorem}
Note that the threshold of $\Omega(\delta k)$ is within a constant factor of optimal, as repeating an algorithm that errs $\delta$ fraction of the time $k$ times will lead to an average of $\delta k$ mistakes.~\Cref{thm:threshold-DPT} implies our standard strong direct sum theorem (\Cref{thm:main-formal}) because \[ \odepthmuk(f^{\otimes k}, \eps, t) \geq \odepthmuk(f^{\otimes k}, \eps, 0) = \odepthmuk(f^{\otimes k}, \eps)\]  for any $t \geq 0$ and $\eps > 0$. 


\subsection{The structure of this section}

By the Hardcore Theorem (\Cref{thm:hardcore theorem for DTs}), proving~\Cref{thm:threshold-DPT} reduces to proving:  \medskip 
%\paragraph{Applying the Hardcore Theorem.} Our proof of \Cref{thm:threshold-DPT} will not directly use $\odepthmu(f, \delta)$. Instead, we prove the following which relates the hardness of $f^{\otimes k}$ to the existence of a hardcore measure for $f$.

\begin{tcolorbox}[colback = white,arc=1mm, boxrule=0.25mm]
\begin{theorem}[Hardness of $f^{\otimes k}$ in terms of a hardcore measure for $f$]
    \label{thm:bounds-from-hardcore}
    Suppose that $f:\bn \to \bits$ has an $(\gamma, d)$-hardcore measure w.r.t $\mu$ of density $\delta$. Then, for any $T:(\bn)^k \to \bits^k$ with $\odepthmuk(T) \leq kd$,%\footnote{We did not attempt to optimize the constants of this theorem.}
    \begin{equation*}
        \Prx_{\bX \sim \mu^k}\bracket*{\lzero{T(\bX) - f^{\otimes k}(\bX)} \leq \lfrac{\delta k}{10}} \leq e^{-\frac{\delta k}{10}} + 10\gamma.
    \end{equation*}
\end{theorem}
\end{tcolorbox}
\medskip 

%Combining \Cref{thm:bounds-from-hardcore} with the hardcore theorem for query complexity, \Cref{thm:hardcore theorem for DTs}, immediately yields \Cref{thm:threshold-DPT}, and so we devote the remainder of 
This section is therefore devoted to proving \Cref{thm:bounds-from-hardcore}.  As discussed in \Cref{sec:tech-overview}, our proof tracks two key quantities: we will analyze how {\sl hardcore density} (\Cref{def:hardcore-density}) and {\sl hardcore advantage} (\Cref{def:hardcore advantage}) are distributed over the leaves of $T$. This proof will be broken into three steps:
\begin{enumerate}
    \item In \Cref{subsec:understand-hardcore}, we prove \Cref{claim:expected total hardcore advantage} and \Cref{lem:resilience lemma}, which aim to understand the distributions of the hardcore density and hardcore advantage of a random leaf of $T$.
    \item In \Cref{subsec:error-one-leaf}, we derive an expression for the probability $T$ makes surprisingly few mistakes as a function of the hardcore density and hardcore advantage at each leaf. This generalizes \Cref{lem:accuracy in terms of hardcore density and advantage}.
    \item In \Cref{subsec:threshold-DPT-proof}, we combine the above  to prove \Cref{thm:bounds-from-hardcore}.
\end{enumerate}

\subsection{How hardcore density and advantage distribute over the leaves}
\label{subsec:understand-hardcore}
We begin with proving our resilience lemma for hardcore density. Roughly speaking, this will say that for any tree $T:(\bn)^k \to \bits^k$, the hardcore density of a random leaf concentrates around $\delta k$. We recall the definition of hardcore density.
\begin{definition}[Hardcore density at $\ell$, \Cref{def:hardcore-density} restated]
\label{def:hardcore-density-body}
For any tree $T:(\bn)^k \to \bits^k$, hardcore measure $H:\bn \to [0,1]$, distribution $\mu$ on $\bn$, $i\in [k]$, and leaf $\ell$ of $T$, the {\sl hardcore density at $\ell$ in the $i$th block} is the quantity: 
\[ \Dens_H(\ell,i)\coloneqq \Ex_{\bX\sim \mu^k}\big[H(\bX^{(i)})\mid \text{$\bX$ reaches $\ell$}\big]. \]
The {\sl total hardcore density at $\ell$} is the quantity $\ds \Dens_H(\ell) \coloneqq \sum_{i =1}^k \Dens_H(\ell,i).$
\end{definition}
The distribution over leaves in the resilience lemma is the canonical distribution.
\begin{definition}[Canonical distribution]
    \label{def:canonical-dist}
    For any tree $T$ and distribution $\mu$ over $T$'s domain, we write $\mu(T)$ to denote the distribution over leaves of $T$ where:  
\[ \Prx_{\bell\sim \mu^k(T)}[\bell = \ell] = \Prx_{\bX\sim \mu}[\text{$\bX$ reaches $\ell$}].\]
\end{definition}
\begin{lemma}[Resilience lemma, generalization of \Cref{lem:resilience lemma}]
    \label{lem:resilience-general}
    For any $T:(\bits^n)^k \to \bits^k$, hardcore measure $H$, distribution $\mu$, and convex $\Phi:\R \to \R$,
    \begin{equation*}
        \Ex_{\bell \sim \mu^k(T)}[\Phi(\dens_H(\bell))] \leq \Ex_{\bz \sim \Bin(k,\delta)}[\Phi(\bz)]
    \end{equation*}
    where $\delta \coloneqq \Ex_{\bx \sim \mu}[H(\bx)]$ is the density of $H$ w.r.t. $\mu$.
\end{lemma}
\begin{proof}
    Draw $\bX \sim \mu^k$. Then, for each $i \in [k]$, independently draw $\bz_i \sim \Ber(H(\bX^{(i)}))$. Note that, for any leaf $\ell$ and $i \in [k]$,
    \begin{equation*}
        \dens_H(\ell, i) \coloneqq \Ex_{\bX \sim \mu^k}[H(\bX^{(i)}) \mid \bX\text{ reaches }\bell] = \Ex_{\bX \sim \mu^k}[\bz_i \mid \bX\text{ reaches }\bell].
    \end{equation*}
    By the above equality and definition $\dens_H(\ell) = \sum_{i \in [k]}\dens_H(\ell,i)$,
    \begin{align*}
        \Ex_{\bell \sim \mu^k(T)}[\Phi(\dens_H(\bell))] &= \Ex_{\bell \sim \mu^k(T)}\bracket[\Bigg]{\Phi\paren[\Bigg]{ \Ex_{\bX \sim \mu^k}\bracket[\Bigg]{\sum_{i \in [k]}\bz_i \mid \bX\text{ reaches }\bell}}}\\
        &\leq \Ex_{\bell \sim \mu^k(T)}\bracket[\Bigg]{\Ex_{\bX \sim \mu^k}\bracket[\Bigg]{\Phi\paren[\Bigg]{ \sum_{i \in [k]}\bz_i \mid \bX\text{ reaches }\bell}}} \tag{Jensen's inequality} \\
        &= \Ex_{\bX \sim \mu^k}\bracket[\Bigg]{\Phi\paren[\Bigg]{ \sum_{i \in [k]}\bz_i}}. \tag{Law of total expectation}
    \end{align*}
    Note that the last line holds precisely for the distribution $\mu^k(T)$ defined in \Cref{def:canonical-dist}, which is why we use that distribution.

    Since $\bX$ is drawn from a product distribution, and $\bz_i$ depends on only the $i^{\text{th}}$ coordinate of $\bX$, $\bz_1,\ldots, \bz_k$ are independent. Furthermore, each has mean $\Ex_{\bx \sim \mu}[H(\bx)] = \delta$. Therefore, $\sum_{i \in [k]}\bz_i$ is distributed according to $\Bin(k,\delta)$.
\end{proof}
A couple of remarks about the above Lemma: First, it implies that $\dens_H(\bell)$ concentrates around $\delta k$. Since $z \mapsto e^{\lambda z}$ is convex for any $\lambda \in \R$,  \Cref{lem:resilience-general} implies that the moment generating function of $\dens_H(\bell)$ is dominated by that of $\Bin(k,\delta)$. This means that Chernoff bounds that hold for $\Bin(k,\delta)$ also hold for $\dens_H(\bell)$. In particular, the statement of \Cref{lem:resilience lemma} is a consequence of the Chernoff bound given in \Cref{fact:chernoff}.

Second, the proof of \Cref{lem:resilience-general} does not make heavy use of the decision tree structure of $T$. It only uses that the leaves of $T$ partition $(\bn)^k$, and so may find uses for other models that partition the domain.


\paragraph{Depth amplification for hardcore advantage.}

While the resilience lemma gives a fairly fine-grained understanding of how hardcore density distributes among the leaves, our guarantee for hardcore advantage are more coarse -- that its expectation over the leaves is bounded.

\begin{definition}[Hardcore advantage at $\ell$, \Cref{def:hardcore advantage} restated]
\label{def:hardcore-advantage-body}
For any tree $T:(\bn)^k \to \bits^k$, hardcore measure $H:\bn \to [0,1]$, distribution $\mu$ on $\bn$, $i\in [k]$, and leaf $\ell$ of $T$, the {\sl hardcore advantage at $\ell$ in the $i$th block} is the quantity: 
\begin{equation}
    \label{eq:hardcore-advantage}
     \Adv_H(\ell,i) \coloneqq \abs[\Big]{\Ex_{\bX\sim \mu^k}\big[ f(\bX^{(i)})T(\bX)_i H(\bX^{(i)})\mid \text{$\bX$ reaches $\ell$}\big]}.
\end{equation}
The {\sl total hardcore advantage at $\ell$} is the quantity  
$ \ds \Adv_H(\ell) \coloneqq \sum_{i=1}^k \Adv_H(\ell, i). $
\end{definition}

\begin{lemma}[Expected total hardcore advantage, \Cref{claim:expected total hardcore advantage} restated]
\label{lem:expected total hardcore advantage body}
    If $H$ is a $(\gamma, d)$-hardcore measure for $f$ of density $\delta$ w.r.t.~$\mu$ and the expected depth of $T$ is at most $dk$, then 
    \begin{equation*}
        \Ex_{\bell \sim \mu^k(T)}[\adv_H(\bell)] \leq \gamma \Ex_{\bell\sim\mu^k(T)}[\Dens_H(\bell)] = \gamma \delta k. 
    \end{equation*}
\end{lemma}
\begin{proof}[Proof of \Cref{lem:expected total hardcore advantage body}]
    By contrapositive. Suppose there exists $T_{\textnormal{large}}:(\bn)^k \to \bits^k$ making $dk$ queries on average w.r.t. $\mu^k$ for which,
    \begin{equation*}
        \Ex_{\bell \sim \mu^k(T)}[\adv_H(\bell)] > \gamma \Ex_{\bell\sim \mu^k(T)}[\dens_H(\bell)] = \gamma\delta k.
    \end{equation*}
    Then, we'll show there exists $T_{\textnormal{small}}:\bn \to \bits$ making $d$ queries on average w.r.t. $\mu$ for which
    \begin{equation*}
        \Ex_{\bx \sim \mu}[ f(\bx) T_{\textnormal{small}}(\bx) H(\bx)] > \gamma \cdot \Ex_{\bx \sim \mu}[H(\bx)] = \gamma \delta.
    \end{equation*}
    Before constructing $T_{\textnormal{small}}$, we observe that we can assume, without loss of generality, that for every leaf $\ell$ of $T$, that we can remove the absolute value from \Cref{eq:hardcore-advantage}; i.e. that
    \begin{equation*}
        \Adv_H(\ell,i) = \Ex_{\bX\sim \mu^k}\big[ f(\bX^{(i)})T(\bX)_i H(\bX^{(i)})\mid \text{$\bX$ reaches $\ell$}\big]
    \end{equation*}
    Otherwise, we could modify this leaf by flipping the label of $T(X)_i$ whenever $X$ reaches a leaf where the above quantity is negative. This does not change the hardcore advantage, so this new $T$ still satisfies our assumption.
    
    $T_{\textnormal{small}}$ will be a randomized algorithm. Upon receiving the input $x \in \bn$, it samples $\bX \sim \mu^k$ and $\bi \sim \mathrm{Unif}([k])$, and then constructs $\bX(x, \bi)$ by inserting $x$ into the $\bi^{\text{th}}$ block of $\bX$,
    \begin{equation*}
        (\bX(x, i))^{(j)} = \begin{cases}
            \bX^{(j)} &\text{if }j \neq i\\
            x&\text{if }j = \bi.
        \end{cases}
    \end{equation*}
    Then, $T_{\textnormal{small}}(x)$ outputs $T_{\textnormal{large}}(\bX(x, \bi))_{\bi}$.

    Our analysis of $T_{\textnormal{small}}$ relies on the following simple observation: If we sample $\bx \sim \mu$, then even conditioning on any choice of $\bi = i$, the distribution of $\bX(\bx, \bi)$ is $\mu^k$. This also means that $\bX(\bx, \bi)$ and $\bi$ are independent.

    We claim that $T_{\textnormal{small}}$ has the two desired properties; low expected number of queries, and high accuracy on $H$. To bound the expected number of queries $T_{\textnormal{small}}$ makes on an input $\bx \sim \mu$, we use that $\bX(\bx, \bi)$ is distributed according to $\mu^k$. Therefore, $T_{\textnormal{large}}(\bX(\bx, \bi))$ makes, on average, $dk$ queries. Expanded, we have that,
    \begin{equation*}
        \sum_{i \in [k], j\in [n]}\Pr[T_{\textnormal{large}}(\bX(\bx, \bi))\text{ queries }\bX(\bx)^{(i)}_j] = dk.
    \end{equation*}
    Whereas, the number of queries $T_{\textnormal{small}}(\bx)$ makes only counts queries to the $\bi^{\text{th}}$ block, and is therefore,
    \begin{align*}
         \sum_{i \in [k], j\in [n]}\Pr\big[T_{\textnormal{large}}(\bX(\bx, \bi))&\text{ queries }\bX(\bx, \bi)^{(i)}_j \cdot \Ind[\bi = i]\big] \\
         &= \sum_{i \in [k], j\in [n]}\Pr\big[T_{\textnormal{large}}(\bX(\bx, \bi)))\text{ queries }\bX(\bx, \bi)^{(i)}_j\big] \cdot \frac{1}{k}\\
         &= d.
    \end{align*}
    In the above, the first equality uses that $\bi$ is independent of $\bX(\bx, \bi)$, and so is still uniform on $[k]$ even conditioned on which queries $T_{\textnormal{large}}$ makes.

    Lastly, we verify that $T_{\textnormal{small}}$ has high accuracy on the hardcore measure.
    \begin{align*}
        \Ex_{\bx \sim \mu}[ f(\bx) T_{\textnormal{small}}(\bx, \bi)) H(\bx)] &= \Ex_{\bx \sim \mu}[ f(\bx) T_{\textnormal{large}}(\bX(\bx, \bi)))_{\bi} H(\bx)] \tag{Definition of $T_{\textnormal{small}}$}  \\
        &=\Ex_{\bi \sim [k]}\bracket*{\Ex_{\bX \sim \mu^k}\bracket*{f(\bX^{(\bi)}) T_{\textnormal{large}}(\bX)_{\bi} H(\bX^{(\bi)})}} \tag{$\bi, \bX(\bx)$ are independent} \\
        &= \Ex_{\bi \sim [k]}\bracket*{\Ex_{\bell \sim \mu^k(T_{\textnormal{large}})}[\adv_H(\bell, \bi)]} \tag{\Cref{def:hardcore-advantage-body}}\\
        &= \frac{1}{k}\Ex_{\bell \sim \mu^k(T_{\textnormal{large}})}[\adv_H(\bell)] > \gamma \delta. \tag*{\qedhere}
    \end{align*}
\end{proof}








\subsection{Understanding the error in terms of hardcore density and advantage}
\label{subsec:error-one-leaf}
To state the main result of this subsection, we'll define the following distribution for the sum of independent Bernoulli random variables.
\begin{definition}
    For any $p \in [0,1]^k$, we use $\BerSum(p)$ to denote the distribution of $\bz \coloneqq \bz_1 + \cdots + \bz_k$ where each $\bz_i$ is independently drawn from $\Ber(p_i)$.
\end{definition}
The following generalizes \Cref{lem:accuracy in terms of hardcore density and advantage}.
\begin{lemma}[Accuracy in terms of hardcore density and advantage of the leaves]
    \label{lem:threshold-accuracy-from-stats}
    Let $H$ be a $(\gamma, d)$-hardcore measure w.r.t.~$\mu$ for $f: \bits^n \to \bits$, and $T:(\bn)^k \to \bits$ be any tree. Then, for any $t \geq 0$, 
    \begin{equation*}
        \Prx_{\bX \sim \mu^k}\bracket*{\lzero[\big]{T(\bX)-  f^{\otimes k}(\bX)} \leq t} \leq \Ex_{\bell \sim \mu^k(T)}\bracket*{\Prx_{\bz \sim \BerSum(p(\bell))}[\bz \leq t]}
    \end{equation*}
    where $p(\ell) \in [0,1]^k$ is the vector where
    \begin{equation*}
        p(\ell)_i \coloneqq \frac{\dens_H(\ell,i) - \adv_H(\ell,i)}{2}\quad\quad\text{for each }i \in [k].
    \end{equation*}
\end{lemma}
\Cref{lem:threshold-accuracy-from-stats} implies a generalization of \Cref{lem:accuracy in terms of hardcore density and advantage}.
\begin{corollary}
    \label{cor:threshold-from-stats-chernoff}
    Let $H$ be a $(\gamma, d)$-hardcore measure w.r.t.~$\mu$ for $f: \bits^n \to \bits$, and $T:(\bn)^k \to \bits$ be any tree. Then, for any $t \geq 0$,
    \begin{equation*}
        \Prx_{\bX \sim \mu^k}\bracket[\Big]{\lzero[\big]{T(\bX)- f^{\otimes k}(\bX)} \leq t} \leq \Ex_{\bell \sim \mu^k(T)}\bracket*{\min\paren[\Big]{1, \exp\paren[\Big]{t - \frac{\Dens_H(\bell) - \Adv_H(\bell)}{4}}}}.
    \end{equation*}
\end{corollary}
\begin{proof}
    The Chernoff bound of \Cref{fact:chernoff} says that, for any $p \in [0,1]^k$ and $\mu \coloneqq \sum_{i \in [k]} p_i$,
    \begin{equation*}
        \Prx_{\bz \sim \BerSum(p)}[\bz \leq t] \leq \begin{cases}
            \exp\paren*{-\frac{(\mu - t)^2}{2\mu}}&\text{if }\mu \geq t \\
            1& \text{otherwise.}
        \end{cases}
    \end{equation*}
    We want to show that the above is bounded by $\min(1, e^{t- \mu/2})$. Clearly this holds for $\mu < t$, so we need only consider the case where $\mu \geq t$
    \begin{equation*}
         \exp\paren*{-\tfrac{(\mu - t)^2}{2\mu}} = \exp\paren*{-\tfrac{\mu^2 - 2t\mu + t^2}{2\mu}} \leq \exp\paren*{-\tfrac{\mu^2 - 2t\mu}{2\mu}} = e^{t - \mu/2}.
    \end{equation*}
    Since $ \exp\paren*{-\frac{(\mu - t)^2}{2\mu}} \leq 1$ as well, it is upper bounded by $\min(1, e^{t- \mu/2})$ as desired. The desired result follows from \Cref{lem:threshold-accuracy-from-stats} as well as $\sum_{i \in [k]} p(\ell)_i = \frac{\Dens_H(\ell) - \Adv_H(\ell)}{2}$ for every leaf $\ell$ of $T$.
\end{proof}



The main observation underlying \Cref{lem:threshold-accuracy-from-stats} is that, if we choose an input $\bX \sim \mu^k$ conditioned on reaching a leaf $\ell \in T$, that $\bX$ is distributed according to a $k$-wise product distribution (i.e.~from $\mu_1(\ell) \times \cdots \times \mu_k(\ell)$ for appropriately defined distributions). The below is essentially the same as  Lemma 3.2 of \cite{Dru12}, but we include a proof for completeness.
\begin{claim}
    \label{claim:leaf-product}
    For any (potentially randomized) tree $T:(\bn)^k \to \mcY$ and leaf $\ell$ of $T$, if $\bX \sim \mu^k$, then the distribution of $\bX$ conditioned on reaching the leaf $\ell$ is a product distribution over the $k$ blocks of $\bX$.
\end{claim}
\begin{proof}
    First, if $T$ is a randomized tree, it is as a distribution over deterministic trees. If the desired result holds for each of those deterministic trees, it also holds for $T$. Therefore, it suffices to consider the case where $T$ is deterministic.

    We'll prove that the distribution of $\bX$ reaching any internal node \emph{or} leaf of $T$ is product by induction on the depth of that node. If that depth is $0$, then all inputs reach it and so the desired result follows from $\mu^k$ being product.

    For depth $d \geq 1$, let $\alpha$ be the parent of $\ell$. Then $\alpha$ has depth $d-1$, so by the inductive hypothesis, the distribution of inputs reaching $\alpha$ is product. Let $i \in [k], j \in [n], b \in \bits$ be chosen so that an input $X$ reaches $\ell$ iff it reaches $\alpha$ and $X^{(i)}_j = b$. Then,
    \begin{align*}
        \Pr[&\bX =X \mid \bX \text{ reaches } \ell] \\
        &= \Pr[\bX =X\mid \bX \text{ reaches } \alpha] \cdot \frac{\Ind[X^{(i)}_j = b]}{\Pr[X^{(i)}_j=b \mid \bX \text{ reaches } \alpha]} \\
         &= \frac{\Ind[X^{(i)}_j = b]}{\Pr[X^{(i)}_j =b\mid \bX \text{ reaches } \alpha]} \cdot \prod_{\ell \in [k]}  \Pr[\bX^{(\ell)} =X^{(\ell)}\mid \bX \text{ reaches } \alpha] \tag{Inductive hypothesis}\\
         &= \paren*{\prod_{\ell \neq i}  \Pr[\bX^{(\ell)} =X^{(\ell)}\mid \bX \text{ reaches } \alpha]} \cdot \frac{\Pr[\bX^{(i)} =X^{(i)}\mid \bX \text{ reaches } \alpha] \cdot \Ind[X^{(i)}_j = b]}{\Pr[X^{(i)}_j =b\mid \bX \text{ reaches } \alpha]}\\
         &= \paren*{\prod_{\ell \neq i}  \Pr[\bX^{(\ell)} =X^{(\ell)}\mid \bX \text{ reaches } \alpha]} \cdot \Pr[\bX^{(i)} =X^{(i)}\mid \bX \text{ reaches } \alpha, \bX^{(i)}_j = b].
    \end{align*}
    The above is decomposed as a product over the $k$ components of $X$, so is a product distribution.
\end{proof}
We conclude this subsection with a proof of \Cref{lem:threshold-accuracy-from-stats}.
\begin{proof}[Proof of \Cref{lem:threshold-accuracy-from-stats}]
    Consider any leaf $\ell$ of $T$. We wish to compute the probability that $T(\bX)$ makes less than $t$ mistakes on $f^{\otimes k}(\bX)$ given that $\bX$ reaches the leaf $\ell$. On this leaf, $T$ outputs a single vector $y \in \bits^k$. Meanwhile, by \Cref{claim:leaf-product}, the distribution of $\bX$ is product over the blocks, and so $f(\bX^{(1)}), \ldots, f(\bX^{(k)})$ are independent. Define $q(\ell) \in [0,1]^k$ as,
    \begin{equation*}
        q(\ell)_i \coloneqq \Prx_{\bX \sim \mu^k}[y_i \neq f(\bX^{(i)})].
    \end{equation*}
    Then,
    \begin{equation*}
        \Prx_{\bX \sim \mu^k}\bracket[\big]{\lzero[\big]{T(\bX), f^{\otimes k}(\bX)} \leq t \mid \bX\text{ reaches }\ell} = \Prx_{\bz \sim \BerSum(q(\ell))}[\bz \leq t].
    \end{equation*}
    For $\bz \sim \BerSum(q)$, the probability $\bz \leq t$ is monotonically decreasing in each $q_i$. Therefore, it suffices to show that $q(\ell)_i \geq p(\ell)_i$ for each $i \in [k]$. We compute,
    \begin{align*}
        q(\ell)_i &= \Prx_{\bX \sim \mu^k}[T(\bX)_i \neq f(\bX^{(i)}) \mid \bX\text{ reaches }\ell] \\
        &=  \frac{1 - \Ex_{\bX \sim \mu^k}[T(\bX)_i f(\bX^{(i)}) \mid \bX\text{ reaches }\ell]}{2} \tag{$T(X)_i, f(X^{(i)} \in \bits$}
    \end{align*}
    Separating the above expectation into two pieces, for $\bX$ drawn from $\mu^k$ conditioned in $\bX$ reaching $\ell$,
    \begin{align*}
        \Ex[T(\bX)_i f(\bX^{(i)})] &= \Ex[T(\bX)_i f(\bX^{(i)})H(\bX^{(i)})] +  \Ex[T(\bX)_i f(\bX^{(i)})(1-H(\bX^{(i)}))] \\
        &\leq \adv_H(\ell,i) +  \Ex[T(\bX)_i f(\bX^{(i)})(1-H(\bX^{(i)}))] \tag{\Cref{def:hardcore advantage}}\\
        &\leq \adv_H(\ell,i) +  \Ex[(1-H(\bX^{(i)}))] = 1 +\adv_H(\ell,i) - \dens_H(\ell,i).
    \end{align*}
    Therefore,
    \begin{equation*}
        q(\ell)_i \geq \frac{\dens_H(\ell,i)- \adv_H(\ell,i)}{2} = p(\ell)_i. \qedhere
    \end{equation*}
\end{proof}

\subsection{Completing the proof of the threshold direct sum theorem}
\label{subsec:threshold-DPT-proof}
In this subsection, we complete the proof of \Cref{thm:bounds-from-hardcore}. Throughout this section, we'll use the following function:
\begin{equation*}
    g_t(z) \coloneqq \min(1, e^{t - z/4}).
\end{equation*}
By using \Cref{cor:threshold-from-stats-chernoff}, it suffices to show that
\begin{equation}
    \label{eq:error-bound}
     \Ex_{\bell \sim \mu^k(T)}\bracket*{g_{\delta k/10}(\Dens_H(\bell) - \Adv_H(\bell))}  \leq e^{-\delta k/10} + 1-\gamma.
\end{equation}
Recall that we have much information about how $\Dens_H(\bell)$ distributes over the leaves via \Cref{lem:resilience-general}, but a coarser understanding of how $\Adv_H(\bell)$ distributes via \Cref{lem:expected total hardcore advantage body}. Because of this, we will first bound the above equation where the $\Adv_H(\bell)$ is set to $0$ and analyze how much including that term affects the result.
\begin{lemma}
    \label{lem:error-no-advantange}
    For any tree $T:(\bn)^k \to \bits^k$ and hardcore measure of density $\delta$ w.r.t. distribution $\mu$,
    \begin{equation*}
        \Ex_{\bell \sim \mu^k(T)}\bracket*{g_{\delta k/10}(\Dens_H(\bell))} \leq e^{-.121\delta k}.
    \end{equation*}
\end{lemma}
\begin{proof}
    We bound,
    \begin{equation*}
        \Ex_{\bell \sim \mu^k(T)}\bracket*{\min\paren[\Big]{1, \exp\paren[\Big]{\delta k / 10 - \frac{\Dens_H(\bell)}{4}}}} \leq e^{\delta k/10} \cdot \Ex_{\bell \sim \mu^k(T)}\bracket*{e^{-\Dens_H(\ell)/4}}.
    \end{equation*}
    Since $z \mapsto e^{-z/4}$ is convex, we can use \Cref{lem:resilience-general} to bound the above using the moment generating function of the binomial distribution,
    \begin{align*}
         \Ex_{\bell \sim \mu^k(T)}\bracket*{e^{-\Dens_H(\bell)/4}} &\leq \Ex_{\bz \sim \Bin(k,\delta)}[e^{-\bz/4}] \\
         &= (1 - \delta(1 - e^{-1/4}))^k \\
         & \leq e^{-(1 -e^{-1/4})\delta k}.
    \end{align*}
    Combining the above,
    \begin{equation*}
         \Ex_{\bell \sim \mu^k(T)}\bracket*{g_{\delta k/10}(\Dens_H(\bell))}  \leq e^{-\delta k (1 - e^{-1/4} - 1/10)} \leq e^{-0.121 \delta k}. \qedhere
    \end{equation*}
\end{proof}
Next, we prove a Lipschitz-style bound for $g$. This will be useful in incorporating $\Adv_H(\bell)$ to \Cref{eq:error-bound}.
\begin{proposition}
    \label{prop:lipschitz-bounds}
    For any $z, \Delta, t \geq 0$,
    \begin{equation}
        \label{eq:lipschitz}
        g_t(z - \Delta) \leq g_t(z) + \Delta/4.
    \end{equation}
    Furthermore, if $z \geq 5t$, then
    \begin{equation}
        \label{eq:lipschitz-scaled}
        g_t(z - \Delta) \leq g_t(z) + \Delta/t.
    \end{equation}
\end{proposition}
\begin{proof}
    \Cref{eq:lipschitz} follows from the $(1/4)$-Lipschitzness of $g_t(z)$.

    For \Cref{eq:lipschitz-scaled}, fix any choice of $z \geq 5t$. We want to show that for any choice of $\Delta$,
    \begin{equation*}
        \frac{ g_t(z - \Delta) - g_t(z)}{\Delta} \leq 1/t.
    \end{equation*}
    We claim that the left hand side of the above inequality is maximized when $\Delta = z - 4t$. When $\Delta$ is increased beyond $z - 4t$, the numerator remains constant (because $g_t(z)$ is constant for any $z \leq 4t$, but the denominator increases, so the maximum cannot occur at at any $\Delta > z-4t$. On the other hand, $g_t(z)$ is convex when restricted to the domain $[4t, \infty)$, so the maximum cannot occur at any $\Delta < z - 4t$. Therefore, it suffices to consider $\Delta = z - 4t$, in which case,
    \begin{equation*}
         \frac{ g_t(z - \Delta) - g_t(z)}{\Delta} = \frac{1 - g_t(z)}{\Delta} \leq \frac{1}{\Delta} \leq 1/t. \qedhere
    \end{equation*}
\end{proof}

We are now ready to prove the main result of this section.
\begin{proof}[Proof of \Cref{thm:bounds-from-hardcore}]
    By applying \Cref{cor:threshold-from-stats-chernoff},
    \begin{equation*}
        \Prx_{\bX \sim \mu^k}\bracket*{\lzero[\big]{T(\bX)-f^{\otimes k}(\bX)} \leq \delta k/10} \leq \Ex_{\bell \sim \mu^k(T)}\bracket*{g_{\delta k/10}(\Dens_H(\bell) - \Adv_H(\bell))}.
    \end{equation*}
    First, we consider the case where $\delta k \leq 40$. Here, by applying \Cref{eq:lipschitz},
    \begin{align*}
        \Ex_{\bell \sim \mu^k(T)}\bracket*{g_{\delta k/10}(\Dens_H(\bell) - \Adv_H(\bell))} &\leq \Ex_{\bell \sim \mu^k(T)}\bracket*{g_{\delta k/10}(\Dens_H(\bell))} + \tfrac{1}{4}\cdot \Ex_{\bell \sim \mu^k(T)}[\Adv_H(\bell)]\\
        &\leq e^{-0.121 \delta k} + \gamma \delta k/4 \tag{\Cref{lem:error-no-advantange,lem:expected total hardcore advantage body}}\\
        &\leq e^{-\delta k/10} + 10\gamma\tag{$\delta k \leq 40$}
    \end{align*}

    When $\delta k > 40$, we break down the desired result into two pieces, depending on whether $\Dens_H(\bell)$ is small or large. For the piece where $\Dens_H(\bell)$ is small, we just use that $g(\cdot)$ is bounded between $0$ and $1$ which means $g(z) - g(z-\Delta) \leq 1$,
    \begin{align*}
         \Ex\big[g_{\delta k/10}(\Dens_H(\bell) &- \Adv_H(\bell)) \cdot \Ind[\Dens_H(\bell) \leq \delta k/2]\big] \\
         &\leq \Ex\bracket*{g_{\delta k/10}(\Dens_H(\bell))\cdot \Ind[\Dens_H(\bell) \leq \delta k/2]} + \Pr[\Dens_H(\bell) \leq \delta k/2]\\
         &\leq  \Ex\bracket*{g_{\delta k/10}(\Dens_H(\bell))\cdot \Ind[\Dens_H(\bell) \leq \delta k/2]}  + e^{-\delta k/8}. \tag{\Cref{lem:resilience lemma}}
    \end{align*}
    For the piece where $\Dens_H(\bell)$ is large, we use \Cref{eq:lipschitz-scaled}
    \begin{align*}
        \Ex\big[g_{\delta k/10}&(\Dens_H(\bell) - \Adv_H(\bell)) \cdot \Ind[\Dens_H(\bell) > \delta k/2]\big] \\
        &\leq \Ex\bracket*{g_{\delta k/10}(\Dens_H(\bell))\cdot \Ind[\Dens_H(\bell) > \delta k/2]} + \tfrac{10}{\delta k} \cdot \Ex\bracket*{\Adv_H(\bell)\cdot \Ind[\Dens_H(\bell) > \delta k/2]} \tag{\Cref{eq:lipschitz-scaled}} \\
        &\leq \Ex\bracket*{g_{\delta k/10}(\Dens_H(\bell))\cdot \Ind[\Dens_H(\bell) > \delta k/2]} + \tfrac{10}{\delta k} \cdot \Ex\bracket*{\Adv_H(\bell)} \tag{$\Dens_H(\bell) \geq 0$} \\
        &\leq \Ex\bracket*{g_{\delta k/10}(\Dens_H(\bell))\cdot \Ind[\Dens_H(\bell) > \delta k/2]} + \tfrac{10}{\delta k} \cdot \gamma \delta k \tag{\Cref{lem:expected total hardcore advantage body}} 
    \end{align*}
    Combining the above two pieces,
    \begin{align*}
        \Ex\big[g_{\delta k/10}(\Dens_H(\bell) - \Adv_H(\bell))\big] &\leq \Ex\big[g_{\delta k/10}(\Dens_H(\bell))\big] + e^{-\delta k/8} + 10\gamma \\
        &\leq e^{-\delta k/8} + e^{-0.121 \delta k} + 10\gamma \tag{\Cref{lem:error-no-advantange}}
    \end{align*}
    When $\delta k > 40$, $e^{-\delta k/8} + e^{-0.121 \delta k} < e^{-\delta k/10}$, so we also recover the desired result in this case.
\end{proof}









